% generated by roam.py handbook export -- не редактировать, править заметки
% 20260923130752-теорема_лагранжа_для_векторнозначных_функции.org (body 0b41906a3917)

\begin{theorem}[{Лагранжа для векторнозначных функций}]\label[theorem]{note:5821D7FE-1F9F-4578-BA2B-D60FDB0C69D3}
  Пусть $E \subset \R^n$~--- открытое \hyperref[note:316B5BFA-8B9F-4267-9086-D6593AAB5F9B]{выпуклое множество}, $g: E \rightarrow \R^m$~--- \hyperref[note:957FA09D-8883-40A5-B4C7-237D9F815A42]{дифференцируемое} отображение.
  Тогда
  \[
  \forall a, b \in E\ \exists c \in [a;b]:\quad \norm{g(b) - g(a)} \le \norm{g^{\prime}(c)} \cdot \norm{b - a}
  \]
  где $g^{\prime}(c)$~--- матрица Якоби $m \times n$, а $\norm{g^{\prime}(c)}$~--- ее спектральная (операторная) норма: $\norm{A} = \sup\limits_{h \ne 0} \frac{\norm{Ah}}{\norm{h}}$.
\end{theorem}

\begin{remark}
  В отличие от одномерного случая, \emph{равенства} $g(b) - g(a) = g^{\prime}(c)(b - a)$ может не быть ни при каком $c$. Например, $g(t) = (\cos t, \sin t)$ на $[0; 2\pi]$: $g(2\pi) - g(0) = 0$, но $g^{\prime}(c) = (-\sin c, \cos c) \ne 0$ при всех $c$.
\end{remark}

\begin{remark}
  Для $g = \nabla f$ матрица Якоби~--- это \hyperref[note:592AD816-1D80-48E6-8399-8524FBC1E8F3]{гессиан} $\nabla^2 f$.
\end{remark}

\begin{proof}
  Пусть $v = g(b) - g(a)$. Если $v = 0$, доказывать нечего. Иначе сведем все к скалярной функции
  \[
  \phi(t) = \left\langle v; g(a + t(b - a)) \right\rangle, \quad t \in [0;1]
  \]
  Она дифференцируема и $\phi^{\prime}(t) = \left\langle v; g^{\prime}(a + t(b-a))(b - a) \right\rangle$.
  По одномерной теореме Лагранжа $\exists \theta \in (0;1)$: $\phi(1) - \phi(0) = \phi^{\prime}(\theta)$. Положим $c = a + \theta(b - a) \in [a;b]$.

  Слева $\phi(1) - \phi(0) = \left\langle v; g(b) - g(a) \right\rangle = \norm{v}^2$. Справа по \hyperref[note:D0601ACD-8C91-4A93-B0E4-AC0D529C6A5C]{КБШ} и определению операторной нормы
  \[
  \left\langle v; g^{\prime}(c)(b - a) \right\rangle \le \norm{v} \cdot \norm{g^{\prime}(c)(b - a)} \le \norm{v} \cdot \norm{g^{\prime}(c)} \cdot \norm{b - a}
  \]
  Итого $\norm{v}^2 \le \norm{v} \cdot \norm{g^{\prime}(c)} \cdot \norm{b - a}$; делим на $\norm{v} > 0$.
\end{proof}
